% Exercises for Chapter 15 (Asian options). Statements only: solutions are kept privately.
% Unless stated otherwise: the test case S = K = 100, T = 1, r = 1%, y = 2%, sigma = 20%,
% constant parameters, fixing dates t_1 < ... < t_m with equal weights 1/m, payment at T.

\begin{exercise}[Means]\label{ex:as-means}
\easy Let $a_1,\dots,a_m>0$ and weights $w_j>0$ with $\sum_jw_j=1$, and let
$H=\big(\sum_jw_j/a_j\big)^{-1}$, $G=\prod_ja_j^{w_j}$ and $A=\sum_jw_ja_j$ be the harmonic,
geometric and arithmetic means.
\begin{enumerate}[label=(\alph*)]
  \item Show that $\min_ja_j\le H\le G\le A\le\max_ja_j$, and that $H=G$, or $G=A$, holds if and
    only if all the $a_j$ are equal.
  \item Order the prices of the calls, and of the puts, with strike $K$ on the harmonic, geometric
    and arithmetic averages of the same fixings.
  \item Let $A$ be the arithmetic average of the fixings $S_{t_j}$, $t_j\le T$, and suppose that
    $\int_t^T(r-y)\ge0$ for every fixing date $t$ (for instance $r\ge y$). Show that the
    average-price call is worth at most the European call with maturity $T$ and the same
    strike. In the test case, where $y>r$, compare the two calls for $m=12$ at $K=60$ and $K=100$.
\end{enumerate}
\end{exercise}

\begin{exercise}[Two fixings]\label{ex:as-two-fixings}
\easy Take $m=2$ fixings at $t_1=0$ and $t_2=T$, with equal weights, so that $A=\frac12(S+S_T)$
and $G=\sqrt{SS_T}$. You may use \noteref{prop:as-geometric}, \noteref{prop:as-moment-matching}
and \noteref{prop:as-curran}.
\begin{enumerate}[label=(\alph*)]
  \item Express the arithmetic call through a European call when $2K>S$, and give its price when
    $2K\le S$. Evaluate both in the test case, for $K=100$ and $K=40$.
  \item Compute $\mu_G$ and $v_G$ from \noteref{prop:as-geometric}(i), once with the double sum
    and once with the tail sums, and price the geometric call for $K=100$.
  \item Show that the three-moment approximation $C_3$ is exact here, and that the lower bound
    $L(x^*)$ equals $C$ for every $K>0$. For which $K$ is the first term $L(\log K)$ exact? Is
    the two-moment approximation exact?
\end{enumerate}
\end{exercise}

\begin{exercise}[A seasoned option]\label{ex:as-seasoned}
\medium A monthly Asian call on the test case ($t_j=j/12$, $K=100$) is valued at $t=\frac12$,
after six fixings, with $S_t=105$. Constant parameters are as in the test case. You may use the
code of the chapter for numerical values.
\begin{enumerate}[label=(\alph*)]
  \item The six known fixings have arithmetic mean $98$. Write the value at $t$ as a multiple of
    the price of an Asian call on the six remaining fixings, valued at $t$, with a shifted strike,
    and compute it.
  \item For the geometric call, the six known fixings have geometric mean $98$. Show that its value
    at $t$ is a Black formula, give its parameters, and compute it.
  \item Now value the arithmetic call at $t=\frac{11}{12}$, just before the last fixing, with
    $S_t=105$ and eleven known fixings of arithmetic mean $100$. Show that it is a multiple of a
    European call, and give its value and delta. Which approximations of the chapter are exact
    in this situation?
\end{enumerate}
\end{exercise}

\begin{exercise}[Thresholds]\label{ex:as-threshold}
\medium Use the notation of \noteref{lem:as-conditional} and \noteref{prop:as-curran}, with
$K>0$ and $v_G>0$, and let $\phi_G$ be the density of $\Lambda=\log G$.
\begin{enumerate}[label=(\alph*)]
  \item Show that $L$ is differentiable with $L'(x)=-D\big(h(x)-K\big)\phi_G(x)$, and deduce
    again that $x^*$ maximises $L$. Why does an error $\delta$ in a finite threshold change $L$
    only by $O(\delta^2)$?
  \item A cheap threshold replaces $h(\log g)$ by $h(\log K)+(g-K)$, its first-order expansion
    around $g=K$ with slope one. Show that the resulting threshold is
    $\hat K=2K-h(\log K)$, and that $L(\log\hat K)\le L(x^*)\le C$ when $\hat K>0$. What lower
    bound remains when $\hat K\le0$? ($L(\log\hat K)$ is the formula published by Curran.)
  \item In the monthly test case ($m=12$, $K=100$) compute $h(\log K)$, $\hat K$, $e^{x^*}$,
    $L(\log K)$, $L(\log\hat K)$ and $L(x^*)$. Repeat for $\sigma=80\%$ and compare with the
    reference price.
\end{enumerate}
\end{exercise}

\begin{exercise}[How good is the geometric control?]\label{ex:as-cv-lab}
\medium This is a computational exercise. Take $K=100$, $T=1$ and equally spaced fixings.
\begin{enumerate}[label=(\alph*)]
  \item For $m\in\{4,12,52\}$ and $\sigma\in\{10\%,20\%,40\%\}$, estimate $1-\rho^2$ for the
    geometric control with $2^{16}$ paths, and compare it with the bound
    $D^2\E[(A-G)^2]/\Var Y$ of \noteref{prop:as-control}, with $\Var Y$ estimated from the same
    paths.
  \item How does $1-\rho^2$ depend on $\sigma$, and how on $m$? Give a heuristic reason for the
    dependence on $\sigma$.
  \item Time the simulation of the paths and the two payoffs, and estimate the efficiency of the
    geometric control relative to plain Monte Carlo for $m=12$, $\sigma=20\%$.
\end{enumerate}
\end{exercise}

\begin{exercise}[An Australian option]\label{ex:as-oz}
\hard An Australian variable-purchase option pays $(A/S_T-k)^+$ at $T$, where $A$ is the average
of the monthly fixings $t_j=jT/12$, $j=1,\dots,12$ (so $t_{12}=T$); constant parameters.
\begin{enumerate}[label=(\alph*)]
  \item Let $X_u=S_{T-u}/S_T$ for $0\le u\le T$. Show that $(X_u)_{0\le u\le T}$ has the law of
    a geometric Brownian motion started at $1$ with volatility $\sigma$ and log-drift $-\nu$, where
    $\nu=r-y-\frac12\sigma^2$. \emph{Hint:} $B_u=W_T-W_{T-u}$ is a Brownian motion.
  \item Deduce that the option is worth an average-price call on $X$ with strike $k$, rate $r$,
    yield $y'=2r-y-\sigma^2$, volatility $\sigma$, fixings at $u_j=T-t_j$ (one of them at $0$)
    and payment at $T$.
  \item For the test case with $k=1$, compute the geometric lower bound, the strip upper bound,
    $C_2$, $C_3$, $L(x^*)$, the conditional lognormal approximation and a reference price, and
    check the reference by simulating $(A/S_T-k)^+$ directly.
  \item A published solution of this exercise uses the yield $2r+y-\sigma^2$. Which price does it
    give, and how would a simulation detect the error?
\end{enumerate}
\end{exercise}

\begin{exercise}[A weekly two-year Asian]\label{ex:as-weekly}
\medium This is a computational exercise. Take weekly fixings $t_j=j/52$, $j=1,\dots,104$,
$T=2$, $\sigma=30\%$, and the other parameters of the test case.
\begin{enumerate}[label=(\alph*)]
  \item Compute a reference price with the recursion \noteref{eq:as-recursion}, and check its
    convergence in the grid size.
  \item Tabulate the errors of $C^G$, $C_2$, $C_3$, the Edgeworth approximation, $L(\log K)$,
    $L(x^*)$ and the conditional lognormal approximation, and compare with
    \noteref{tab:as-approx}.
  \item With $2^{16}$ paths, compare plain Monte Carlo, the geometric control, the three controls
    and the conditional estimator of \noteref{prop:as-rest}: standard deviations per path and
    estimates with their standard errors.
\end{enumerate}
\end{exercise}
